\begin{document}

\section{関連研究と課題設定}

\subsection{検知後調査としての行動列復元}

エンドポイント調査では，Windows Securityログ，Sysmon，EDR telemetry，DNS，ブラウザ履歴など，粒度や記録目的の異なるログを横断する．Sysmonはプロセス生成，親プロセス，コマンドライン，ネットワーク接続，レジストリ操作などを記録できる\cite{sysmon}．EDRは検知理由やプロセス情報をまとめたalert summaryを提供する場合がある．これらは調査に有用だが，alert summaryは観測された全行動を代表するものではない．したがって，最終出力はアラート文面ではなく，ログ上で確認できる行動主張に限定する必要がある．

本研究の課題は，異常検知モデルの精度評価ではない．DeepLogのような系列異常検知研究は，ログ系列から異常を検出・診断することに焦点を置く\cite{deeplog}．一方，本研究は検知後のSOC調査において，調査者が検証可能な証跡付き行動列を復元できるかを問う．この違いは評価単位にも表れる．本研究では自然文の類似度ではなく，主体，操作，対象，コマンドライン，証跡行などの必須itemを採点する．

\subsection{LLMエージェントによる探索}

LLMエージェントは，初期手がかりから仮説を立て，外部環境への問い合わせ結果に応じて次の探索を変更できる．ReActは推論と行動を交互に進める枠組みを示しており\cite{react}，ログ調査でも，時刻範囲，プロセス木，ファイル・レジストリ対象を段階的に確認する流れと親和性がある．CLOUSEAUは，raw incident logsをLLMが扱いやすい構造化形式へ変換し，QA agentを通じてログを検索可能にし，Investigator agentとChief Inspector agentが調査leadの生成，証跡収集，attack storyの統合を行う\cite{clouseau}．本研究はこの設計思想に基づくが，評価対象を「攻撃story全体の復元」から「SOCアラートやprocess/time clue周辺の良性を含む行動chainの証跡付き復元」へ移す．

\subsection{ATLASv2の利用範囲}

ATLASv2は，攻撃エンゲージメントと通常業務活動を含むログデータセットである\cite{atlasv2}．本研究では，ATLASv2を「攻撃検知済みの答えを読むデータ」ではなく，「SOCでアラートやtelemetryを受け取った後，端末上の行動を証跡付きで説明する能力を測るデータ」として利用する．対象は，DNS packet capture，Python SimpleHTTPServer，Sublime経由のscript execution，cmd.exe周辺操作，Discord Run Key関連操作など，良性環境内でもセキュリティアラートや調査対象になり得る行動である．

\end{document}
